\documentclass[10pt,a4paper]{article}

\usepackage[hmargin=2cm,vmargin=1.8cm]{geometry}
\usepackage[T1]{fontenc}
\usepackage[utf8]{inputenc}
\usepackage{lmodern}
\usepackage{microtype}
\usepackage{booktabs}
\usepackage{amsmath}
\usepackage[numbers,sort&compress]{natbib}
\usepackage[hidelinks]{hyperref}

\title{\vspace{-1.5em}\textbf{Knowledge Editing via Self-Study Data Synthesis}\\[0.3em]
\large Abschlussbericht Projektarbeit}
\author{Iraj Masoudian, Heinrich-Heine-Universit\"at D\"usseldorf} % TODO: supervisor / chair line if required
\date{17.07.2026}

\setlength{\abovecaptionskip}{4pt}
\setlength{\belowcaptionskip}{0pt}

\begin{document}
\maketitle
\vspace{-3em}

\section{Methodology}
The goal of this project is to inject factual edits into a language model by training on
synthetic conversations generated with the Self-Study procedure of the Cartridges
framework \citep{eyuboglu2025cartridgeslightweightgeneralpurposelong}, and to compare this approach with dedicated
knowledge editing methods such as ROME \citep{meng2023locatingeditingfactualassociations}, MEMIT \citep{meng2023masseditingmemorytransformer},
and AnyEdit \citep{jiang2025anyediteditknowledgeencoded}.

\subsection{Benchmark}
All experiments use the CounterFact split of the AKEW benchmark \citep{wu2024akewassessingknowledgeediting} with
975 counterfactual edits. Each entry provides the subject, the old (true) answer, the new
(edited) answer, one base edit question, two paraphrases, and unstructured sentence
representations of the fact.

\subsection{Data Synthesis with Self-Study}
Self-Study is a loop between two agents: a user bot asks questions about a given context,
in our setting a single CounterFact edit, and an assistant bot answers them. Both are
instances of the teacher model (\texttt{Qwen3-4b}, served with the Tokasaurus inference
server). Every conversation starts from a seed prompt, an instruction from which the user
bot derives its first message.

The data sources shipped with Cartridges handle plain text corpora. A new resource for
knowledge editing was therefore implemented: it loads the CounterFact entries, extracts
the subject, the relation, and the old and new answer of each edit, derives the relation
phrase automatically from the CounterFact prompt template, and passes the unstructured
fact sentence to the synthesizer as the conversation context.

The seed prompts that ship with the Cartridges framework target general corpora and
cover generic tasks such as structuring, summarization, question answering, use cases,
and creative writing; they are not suited for targeted fact edits. As part of this
project, new seed prompt types were therefore designed specifically for knowledge
editing. Each is instantiated with the subject, the relation, and the old or new answer
of the current edit:
\begin{itemize}\itemsep2pt
  \item \emph{Negation}: the user asks whether the old fact still holds or asserts it
        with varying confidence; the assistant denies it and states the new fact.
  \item \emph{Correction}: the user insists on the old fact and is corrected.
  \item \emph{Question}: the user asks directly for the edited relation, in varied
        phrasings and registers.
  \item \emph{Strict blocking}: the user probes the former connection between the
        subject and the old answer.
  \item \emph{Creative}: an adapted variant of the generic creative type that explores
        consequences of the edit.
\end{itemize}
The pipeline itself was modified in several points. Each seed type now injects its own
system prompt, which shapes the assistant's behavior for that conversational pattern and
instructs it to present the edit as its own knowledge; originally all seed prompts
shared one generic system prompt. The seed prompt is also no longer passed to the user
bot verbatim as its first message: it only describes an intent, which the user bot
rewrites into a natural message in the voice of a randomly assigned persona, sampled at
a high temperature of $0.95$ for lexical diversity. The assistant answers at a low
temperature of $0.4$; assistant thinking is enabled with a fixed probability. With this setup, 4{,}096 samples per seed type were generated,
16{,}384 in total, so every edit appears about 17 times (4 to 5 samples per seed type
and entry). In addition, single entry datasets were created to train per edit cartridges
that can be concatenated.

\subsection{Training}
Two variants were trained on the synthesized data. The cartridge setup consists of one
large cartridge containing all 975 edits and several small cartridges containing one edit
each, which can be composed at inference time. Cartridges are initialized from the KV
cache of a prompt; a new initializer was written that builds this prompt directly from a
CounterFact entry as a set of behavior rules (state the new answer, correct mentions of
the old one, do not reconnect the subject with the old answer) instead of generic corpus
text. As a comparison, a LoRA adapter
\citep{hu2021loralowrankadaptationlarge} was trained on the same data (configuration in Table~\ref{tab:lora}).

\begin{table}[ht]
\centering
\footnotesize
\begin{tabular}{llll}
\toprule
\multicolumn{2}{c}{\textbf{Model \& LoRA}} & \multicolumn{2}{c}{\textbf{Training}}\\
\cmidrule(lr){1-2}\cmidrule(lr){3-4}
Base model    & Qwen3-4b & Batch size            & 4\\
LoRA $r$      & 16       & Grad.\ accumulation   & 8\\
LoRA $\alpha$ & 32       & Effective batch size  & 32\\
LoRA dropout  & 0.05     & Epochs                & 5\\
Max length    & 2048     & Learning rate         & $2\cdot10^{-4}$\\
Precision     & BF16     & Warmup steps          & 100\\
\midrule
Train samples & 14{,}745 & Eval samples          & 1{,}639\\
\bottomrule
\end{tabular}
\caption{LoRA training configuration of the interim results.}
\label{tab:lora}
\end{table}

\subsection{Evaluation}
Evaluation follows two tracks. First, an LLM judge scores model responses from 0 to 5
along four standard dimensions: efficacy (original question, is the edited fact
returned), generalization (paraphrased question), locality (does the edit leak into
unrelated facts), and portability (can the model reason on top of the edit). Second,
BERTScore and ROUGE-L are computed against the reference answers, separately for original
and paraphrased questions, matching the AKEW evaluation protocol used by
AnyEdit \citep{jiang2025anyediteditknowledgeencoded}.

\section{Interim Results}
Table~\ref{tab:results} shows both variants next to the published AKEW CounterFact
results for Qwen2.5-7B-Instruct.

\begin{table}[ht]
\centering
\footnotesize
\begin{tabular}{lcccc}
\toprule
& \multicolumn{2}{c}{\textbf{Original}} & \multicolumn{2}{c}{\textbf{Paraphrased}}\\
\cmidrule(lr){2-3}\cmidrule(lr){4-5}
\textbf{Method} & BERTScore & ROUGE-L & BERTScore & ROUGE-L\\
\midrule
Pre-edited                        & 65.50\,$\pm$\,0.34 & 18.24\,$\pm$\,0.60 & 44.74\,$\pm$\,0.41 & 17.29\,$\pm$\,0.51\\
FT-L                              & 46.66\,$\pm$\,0.48 & 14.63\,$\pm$\,0.58 & 32.34\,$\pm$\,0.50 & 12.31\,$\pm$\,0.62\\
MEND                              & 69.54\,$\pm$\,0.54 & 25.47\,$\pm$\,0.49 & 52.86\,$\pm$\,0.40 & 22.81\,$\pm$\,0.54\\
ROME~\citep{meng2023locatingeditingfactualassociations}         & 75.89\,$\pm$\,0.38 & 36.42\,$\pm$\,0.45 & 55.67\,$\pm$\,0.47 & 25.79\,$\pm$\,0.59\\
MEMIT~\citep{meng2023masseditingmemorytransformer}       & 77.19\,$\pm$\,0.32 & 38.95\,$\pm$\,0.48 & 56.04\,$\pm$\,0.40 & 25.73\,$\pm$\,0.57\\
AlphaEdit                         & 80.66\,$\pm$\,0.25 & 45.55\,$\pm$\,0.37 & 56.99\,$\pm$\,0.49 & 27.69\,$\pm$\,0.59\\
AnyEdit~\citep{jiang2025anyediteditknowledgeencoded}  & 98.08\,$\pm$\,0.15 & 95.09\,$\pm$\,0.19 & 65.40\,$\pm$\,0.38 & 43.49\,$\pm$\,0.47\\
UnKE                              & 97.34\,$\pm$\,0.13 & 90.44\,$\pm$\,0.16 & 59.29\,$\pm$\,0.48 & 29.27\,$\pm$\,0.61\\
AnyEdit$^{*}$                     & 99.63\,$\pm$\,0.09 & 98.99\,$\pm$\,0.10 & 60.78\,$\pm$\,0.39 & 32.95\,$\pm$\,0.59\\
\midrule
LoRA (ours)                       & 87.90 & 29.71 & 86.35 & 22.45\\
Cartridge (ours)                  & 90.15 & 44.92 & 87.56 & 27.31\\
\bottomrule
\end{tabular}
\caption{AKEW CounterFact results (BERTScore / ROUGE-L, original and paraphrased
questions). Baseline rows are the numbers published in~\citep{jiang2025anyediteditknowledgeencoded} for
Qwen2.5-7B-Instruct; the last two rows are interim measurements from this project.}
\label{tab:results}
\end{table}

Both variants reach higher BERTScore values than the locate and edit baselines (MEND,
ROME, MEMIT, AlphaEdit), and the cartridge comes close to AlphaEdit on the ROUGE-L of the
original questions. The cartridge performs better than the LoRA adapter trained on the
same data in all four columns. The difference between original and paraphrased questions
is small for both variants (cartridge: 90.15 vs.\ 87.56 BERTScore), while the baselines
lose considerably more under paraphrasing. AnyEdit remains well ahead on the original
questions. Two caveats apply to these numbers: the LLM judge frequently produced
unparseable JSON output, so judge scores are not reported here, and both variants
finetune \texttt{Qwen3-4b} while the baseline rows are published for Qwen2.5-7B-Instruct,
so the comparison is only indicative.

\section{Bugfixes and Optimizations}
After the interim results, the pipeline was revised in the following points.

\begin{itemize}\itemsep2pt
  \item \textbf{Judge output.} The judge now runs on a vLLM server with grammar
        constrained decoding (xgrammar), which guarantees valid JSON. The parsing
        failures that previously invalidated a large share of the scores are eliminated,
        and every evaluation run is gated by a calibration test with known answers.
  \item \textbf{Training data contamination.} An audit of the synthesized dataset showed
        two defects: raw \texttt{<think>} blocks stored as training targets (about 20
        percent of the samples) and answers that reference their source (``according to
        the provided information'', about 25 percent). Thinking is now disabled during
        synthesis, the negation and question system prompts were rewritten so that the
        assistant states facts directly without citing a source, and a filter removes
        remaining violations after synthesis.
  \item \textbf{Teacher fidelity.} A second audit showed that in some answers the
        teacher keeps the old fact, often by presenting old and new fact as compatible.
        Every sample is therefore checked by the LLM judge, and answers that still
        assert the old fact are removed (13.3 percent of the regenerated dataset).
  \item \textbf{Edit coverage.} Edits are now drawn round robin instead of at random, so
        all 975 edits are covered uniformly, and every sample records the edit it
        belongs to.
  \item \textbf{Metric correction.} The check for mentions of the old fact used a plain
        substring search and counted, among others, matches inside subject names. It now
        matches on word boundaries and skips locality tests, where the old answer is
        often the correct one.
  \item \textbf{Student model.} The student was switched to Qwen2.5-7B-Instruct, the
        model the published baselines edit. A new dataset of 36{,}864 samples was
        synthesized with the corrected pipeline; after both filter stages 30{,}442
        samples remain. The LoRA adapter was retrained on this dataset
        ($r=64$, $\alpha=256$, learning rate $2\cdot10^{-5}$, 5 epochs).
\end{itemize}

The retrained adapter was evaluated on all 975 CounterFact edits. Table~\ref{tab:akewnew}
reports the AKEW protocol metrics in the format of Table~\ref{tab:results}, so they can
be compared directly with the interim results and the published baselines.

\begin{table}[ht]
\centering
\footnotesize
\begin{tabular}{lcccc}
\toprule
& \multicolumn{2}{c}{\textbf{Original}} & \multicolumn{2}{c}{\textbf{Paraphrased}}\\
\cmidrule(lr){2-3}\cmidrule(lr){4-5}
\textbf{Method} & BERTScore & ROUGE-L & BERTScore & ROUGE-L\\
\midrule
LoRA (ours, retrained) & 91.26 & 47.28 & 88.38 & 31.97\\
\bottomrule
\end{tabular}
\caption{AKEW CounterFact results of the retrained LoRA adapter (Qwen2.5-7B-Instruct,
cleaned dataset) on all 975 edits, in the format of Table~\ref{tab:results}.}
\label{tab:akewnew}
\end{table}

Table~\ref{tab:final} breaks the evaluation down by test type, 6{,}825 test cases in
total: per edit one efficacy question, two paraphrases, two locality probes drawn from
the neighborhood and attribute prompts, and two portability questions drawn from the
generation prompts.

\begin{table}[ht]
\centering
\footnotesize
\begin{tabular}{lcccccc}
\toprule
\textbf{Test type} & BERTScore & ROUGE-L & Exact match & Judge (0--5) & Success & Old fact leak\\
\midrule
Efficacy       & 91.26 & 47.28 & 83.4\,\% & 4.44 & 85.4\,\% & 14.4\,\%\\
Generalization & 88.38 & 31.97 & 54.3\,\% & 3.98 & 72.6\,\% & 26.0\,\%\\
Locality       & 85.26 & 19.85 & 22.2\,\% & 4.84 & 96.0\,\% & --\\
Portability    & 88.51 & 28.26 & 56.0\,\% & 3.54 & 62.6\,\% & 20.7\,\%\\
\midrule
Overall        & 87.94 & 29.64 & 49.7\,\% & 4.17 & 78.3\,\% & 15.4\,\%\\
\bottomrule
\end{tabular}
\caption{Evaluation of the retrained LoRA adapter (Qwen2.5-7B-Instruct, cleaned dataset)
on all 975 CounterFact edits (975 efficacy tests, 1{,}950 tests for each other type).
Success is the share of judge scores of 4 or higher; the old fact check is skipped for
locality tests, where the old answer is the correct one.}
\label{tab:final}
\end{table}

The edits themselves are learned reliably: on the original questions the adapter reaches
a judge score of 4.44 and answers 83.4 percent of them with the exact new target.
Locality is largely preserved at 96 percent success, so the training rarely damages
unrelated facts. Generalization and portability are weaker, and in about a quarter of
the paraphrase answers the old fact still appears. On the AKEW protocol
(Table~\ref{tab:akewnew}) the adapter improves on the interim adapter in all four
values. Since the student is now Qwen2.5-7B-Instruct, these numbers are also directly
comparable to Table~\ref{tab:results}: the adapter surpasses all locate and edit baselines on
BERTScore, exceeds AlphaEdit on the original ROUGE-L, and keeps its small gap between
original and paraphrased questions, while AnyEdit remains ahead on the original
questions.

\begingroup
\footnotesize
\setlength{\bibsep}{2pt}
\bibliographystyle{unsrtnat}
\bibliography{references}
\endgroup

\end{document}
